\documentclass[tikz,border=18pt]{standalone}
\usepackage{luatexja}
\usepackage{tikz}
\usetikzlibrary{shapes.geometric,arrows.meta,positioning,calc}

\definecolor{userblue}{RGB}{20, 80, 150}
\definecolor{badred}{RGB}{190, 30, 30}
\definecolor{goodgreen}{RGB}{25, 125, 45}

\begin{document}
\begin{tikzpicture}[>=Stealth, font=\small]

  % -------------------------------------------------------------
  % 上段：公衆Wi-Fiでの暗号化なし通信（中間者攻撃・盗聴）
  % -------------------------------------------------------------
  \node[draw=badred, fill=red!6, rounded corners=8pt, line width=1.5pt, minimum width=21.0cm, minimum height=4.2cm, align=left] (unsec_box) at (0, 4.0) {};
  
  \node[above right, font=\large\bfseries, text=badred] at (-10.2, 5.3) {【危険】暗号化なし公衆Wi-Fi（HTTP平文通信・なりすましAP）};

  \node[draw=userblue, fill=blue!10, rounded corners=6pt, thick, align=center, minimum width=3.8cm, minimum height=1.8cm] (user1) at (-7.8, 3.4) {
    \textbf{\large 高校生スマホ}\\[3pt]
    \footnotesize (カフェで利用)
  };

  \node[draw=badred, fill=red!15, rounded corners=6pt, very thick, align=center, minimum width=4.6cm, minimum height=1.8cm] (hacker) at (0, 3.4) {
    \textbf{\large \textcolor{badred}{攻撃者（中間者: MitM）}}\\[3pt]
    \footnotesize 同一Wi-Fi上の盗聴・パケット解析
  };

  \node[draw=black!70, fill=gray!15, rounded corners=6pt, thick, align=center, minimum width=3.8cm, minimum height=1.8cm] (server1) at (7.8, 3.4) {
    \textbf{\large Webサーバー}\\[3pt]
    \footnotesize (HTTPサービス)
  };

  % 矢印とラベル：ノードから十分離れた位置に独立配置
  \draw[->, thick, dashed, badred] (user1.east) -- (hacker.west)
    node[midway, above=4pt, font=\footnotesize\bfseries, text=badred] {平文データ};
  \draw[->, thick, dashed, badred] (user1.east) -- (hacker.west)
    node[midway, below=4pt, font=\scriptsize, text=badred] {(パスワード/閲覧履歴)};

  \draw[->, thick, dashed, badred] (hacker.east) -- (server1.west)
    node[midway, above=4pt, font=\footnotesize\bfseries, text=badred] {覗き見・改ざん};

  % -------------------------------------------------------------
  % 下段：HTTPS / VPN による暗号化トンネル保護
  % -------------------------------------------------------------
  \node[draw=goodgreen, fill=green!6, rounded corners=8pt, line width=1.5pt, minimum width=21.0cm, minimum height=4.2cm, align=left] (sec_box) at (0, -1.6) {};
  
  \node[above right, font=\large\bfseries, text=goodgreen] at (-10.2, -0.3) {【安全】HTTPS (TLS暗号化) ＆ VPN による強固なシールドトンネル};

  \node[draw=userblue, fill=blue!10, rounded corners=6pt, thick, align=center, minimum width=3.8cm, minimum height=1.8cm] (user2) at (-7.8, -2.2) {
    \textbf{\large 高校生スマホ}\\[3pt]
    \footnotesize (暗号化通信)
  };

  \node[draw=goodgreen, fill=green!20, rounded corners=6pt, very thick, align=center, minimum width=8.2cm, minimum height=1.8cm] (tunnel) at (0, -2.2) {
    \textbf{\large \textcolor{goodgreen}{TLS / VPN 暗号化トンネル}}\\[4pt]
    \small \textbf{共通鍵暗号(AES)＋公開鍵暗号(RSA/ECC)のハイブリッド}\\[2pt]
    \footnotesize ※通信内容は完全暗号化され、第三者には解読不能！
  };

  \node[draw=goodgreen, fill=green!15, rounded corners=6pt, thick, align=center, minimum width=3.8cm, minimum height=1.8cm] (server2) at (7.8, -2.2) {
    \textbf{\large 正規Webサーバー}\\[3pt]
    \footnotesize (HTTPSサービス)
  };

  \draw[->, very thick, goodgreen] (user2.east) -- (tunnel.west);
  \draw[->, very thick, goodgreen] (tunnel.east) -- (server2.west);

\end{tikzpicture}
\end{document}
